\section{Recovering data through hidden-opponent tests}
Here actions are committed, or tested as singleton menus. There are $m$
opponent positions with a common alphabet of $k$ actions. Hiding a pair
$e=\{r,s\}$ averages a discrepancy array $d$ by
$E_e d=(d+d\circ\tau_e)/2$. Other positions remain distinguished.
Varying product laws and a common IID law at the hidden pair recovers this
averaged array by polynomial coefficients.

\begin{theorem}[Pair-mask criterion; P7-O1--O2]
Let the tested pairs form a graph whose component sizes are $r_1,\ldots,r_q$.
The space invisible to every pair test has dimension
\[
 \prod_{j=1}^q\binom{k}{r_j}.
\]
Thus ordered data are recovered exactly when some component has more than
$k$ positions. If $m>k$, exactly $k$ pair contexts are necessary and sufficient
for uniform recovery. If $m\le k$, even every pair fails to recover all data.
\end{theorem}
\begin{proof}
Invisibility means $d\circ\tau_e=-d$ for every edge. Edge transpositions
generate the full symmetric group on each connected component, so $d$ is
alternating within each component. Such arrays vanish on repeated labels
and are freely determined by one value per subset of $r_j$ distinct labels.
Tensoring components gives the formula. A component on $k+1$ vertices
requires at least $k$ edges, and a path with $k$ edges attains the bound.
\end{proof}
Applying the criterion simultaneously to every recipient and focal action
recovers an exact typed automorphism only when the \emph{same} typed map is
used throughout. Independently chosen focal maps do not suffice.

For unoptimized payoff expectations, all common IID tests are equivalent to
all exchangeable tests: coefficients of the homogeneous IID polynomial
recover every histogram average (P7-O3). The optimized example above shows
why the quantifier over actions cannot be moved through this comparison.

With common ordered physical consequences, indifference under every strictly
increasing utility identifies the entire hidden consequence distribution by
its tail probabilities. A hidden pair therefore identifies its unordered
pair of consequences. Apply the alternating-array argument to each tail
indicator to obtain pointwise consequence recovery under the same graph
condition (P7-O4). The proofs and sharp alternating counterexamples are in
\texttt{research/phase7-independent-models.md}. Alternating tensors and
polynomial identification are standard tools; novelty of this particular
experimental application remains unestablished.
